\chapter{平面圆锥曲线}

\section{射影空间}

在复分析中，为了将复平面紧化，引入了添加上无穷远点后的复平面$\mathbb{C}\cup\{\infty\}$，即\bluebf{扩充复平面}，并且扩充复平面上的任意一个点和一个复球面上的点一一对应.\par
现在，让我们从另一个视角来看扩充复平面.将复数域零点去掉后，$\mathbb{C}^*=\mathbb{C}\setminus\{0\}$对复数乘法成群，因此它对$\mathbb{C}^2$有如下的群作用：$\setjot[-8]\begin{aligned}
    \mathbb{C}^* \times \mathbb{C}^2 &\to \mathbb{C}^2 \\
    (\lambda,(x,y)) &\mapsto (\lambda x,\lambda y)
\end{aligned}$.\par
下一步，在$\mathbb{C}^2\setminus\{0\}$上商掉$\mathbb{C}^*$，定义$\mathbb{C}\mathbb{P}^1=\left(\mathbb{C}^2\setminus\{0\}\right)/\mathbb{C}^* = \left\{ [x_1,x_2] | (x_1,x_2)\in \mathbb{C}^2\setminus\{0\}\right\}$，称为\bluebf{射影空间}，其中的等价关系可以描述为：若$x_2\neq 0,(x_1,x_2)\sim \left(\frac{x_1}{x_2},1\right)$.\par


\begin{definition}{$n$维射影空间}
    定义\textbf{$n$维射影空间}为$\mathbb{C}\mathbb{P}^n=\left(\mathbb{C}^{n+1}\setminus\{0\}\right)/\mathbb{C}^*$，其中，$\mathbb{C}^*$对$\mathbb{C}^{n+1}$的群作用为：
    \begin{align*}
        \mathbb{C}^*\times \mathbb{C}^{n+1} &\to \mathbb{C}^{n+1} \\
        (\lambda,(x_1,x_2,\cdots,x_n)) &\mapsto (\lambda x_1,\lambda x_2,\cdots,\lambda x_n)
    \end{align*}
\end{definition}
\begin{remark}
    $\mathbb{C}\mathbb{P}^n$为紧集.
\end{remark}



\begin{wrapfigure}{r}{0.3\textwidth}
    \vspace{-1 em}
    \begin{tikzcd}[every label/.append style={font=\small},font=\small]
    & & {\begin{gathered}
        \mathbb{C}^*=\mathbb{C}\setminus\{0\}\\
        \dfrac{x_2}{x_1}=z
    \end{gathered}}
    \arrow[dd,"\varphi"] \\
    {\begin{gathered} U_1\cup U_2\\ [x_1,x_2]\end{gathered}}
    \arrow[rru,"\varphi_1"] \arrow[drr,"\varphi_2"'] & & \\
    & & {\begin{gathered}\mathbb{C}^*=\mathbb{C}\setminus\{0\}\\\dfrac{x_1}{x_2}=\dfrac{1}{z}\end{gathered}}
\end{tikzcd}
    \captionof{figure}{$\mathbb{C}\mathbb{P}^1$和$\mathbb{C}^*$的关系}
    \vspace{-5 em}
\end{wrapfigure}
下面介绍$\mathbb{C}\mathbb{P}^1$的性质与上面的某些结构.\par
首先，定义$U_1=\{[x_1,x_2]|x_1\neq 0\}$和$U_2=\{[x_1,x_2]|x_2\neq 0\}$，则$\mathbb{C}\mathbb{P}^1=U_1\cup U_2$.\par
接下来，我们就可以在$U_1$上定义到$\mathbb{C}$的映射：$\setjot[-6]\begin{aligned}[t]
    \varphi_1: U_1 &\xrightarrow{\cong} \mathbb{C} \\
    [x_1,x_2] &\mapsto \frac{x_2}{x_1}    
\end{aligned}$.\par
对上述映射，同时可以定义其逆映射$\setjot[-6]\begin{aligned}[t]
    \varphi_1^{-1}: \mathbb{C} &\rightarrow U_1 \\
    x &\mapsto [1,x]
\end{aligned}$.\par
因此我们找到了$U_1$和$\mathbb{C}$之间的双射，所以$U_1$和$\mathbb{C}$是相同的.\par
同样的道理可以定义$U_1$到$\mathbb{C}$之间的同构.\par
综上所述，$\mathbb{C}\mathbb{P}^1=U_1\cup U_2$和$\mathbb{C}^*$之间可以形成如右图所示的交换图.

因此，现在我们就有两种方式来理解$\mathbb{C}\mathbb{R}^1$.其一是扩充复平面$\mathbb{C}\cup\{\infty\}$，其二是$\mathbb{C}^2$的商空间$(\mathbb{C}-\{0,0\})/\mathbb{C}^*$.其中第二种方法$\mathbb{C}^2$可以推广到任意域上的有限维线性空间$V:k(\mathbb{F}_p,\mathbb{Q},\mathbb{R})$，定义$\mathbb{P}(V)=(V\setminus\{0\})/k^*$.例如可以定义$\mathbb{C}\mathbb{P}^2=\left(\mathbb{C}^3\setminus\{0\}\right)/\mathbb{C}^*=\mathbb{C}^2\cup(\mathbb{C}\mathbb{P}^1)$.

下面一个话题是复分析中的分式线性变换.一个分式线性变换可以写成
\begin{align*}
    f: \mathbb{C}\cup\{\infty\} &\to \mathbb{C}\cup\{\infty\} \\
    z &\mapsto \frac{az+b}{cz+d},\quad \begin{pmatrix}
        a & b \\
        c & d
    \end{pmatrix}\text{是非退化二阶方阵}
\end{align*}
如果从射影空间的角度来看，上面的映射可以写成
\begin{align*}
    f: \mathbb{C}\mathbb{P}^1 &\to \mathbb{C}\mathbb{P}^1 \\
    [x_1,x_2]=\left[1,\frac{x_2}{x_1}\right] &\mapsto \left[1,\frac{a\tfrac{x_2}{x_1}+b}{c\tfrac{x_2}{x_1}+d}\right] = [cx_2+dx_1,ax_2+bx_1]
\end{align*}
